% !TeX root = main.tex
\chapter{Introduction}
\label{ch:introduction}

Le clustering $k$-means classique suppose que l'ensemble des points tient en
mémoire et peut être relu autant de fois que nécessaire. Cette hypothèse
s'effondre dès que les données arrivent en flux continu : un capteur, un
flux de clics, un flux de transactions. On ne peut ni tout stocker, ni
relire deux fois le même point, et pourtant on veut un clustering de qualité
comparable à ce qu'on obtiendrait en traitant les données par lots.

\textsc{StreamKM++} (Ackermann et al., 2012) répond à ce problème par la
technique du \emph{coreset} : au lieu de conserver le flux, on en maintient
un résumé pondéré de taille bornée $m$, construit incrémentalement, tel que
le coût de clustering calculé sur ce résumé approxime le coût de clustering
sur les données complètes, pour n'importe quel jeu de centres (Définition 2.2
de la thèse de référence). Un seul passage sur les données suffit, et la
mémoire utilisée ne dépend que de $m$ et $\log(n)$, jamais de $n$
directement.

\paragraph{Objet de ce rapport.} Bitsakis (2018) implémente StreamKM++ dans
Apache Flink et l'évalue de façon approfondie (chapitres 4 et 5 de sa
thèse). Notre travail consiste à \textbf{porter cette implémentation vers
Apache Spark}, dans le cadre du cours de déploiement ML à grande échelle de
ce M2. Ce portage n'est pas une simple traduction syntaxique : Flink et
Spark reposent sur des modèles d'exécution différents (état par
\emph{subtask} adressable vs. RDD immuables et micro-batch), et une partie
du travail intéressant de ce rapport consiste précisément à identifier ce
qui se traduit directement, ce qui demande une réinterprétation, et ce qui
n'a tout simplement pas d'équivalent côté Spark (chapitre~\ref{ch:implementation},
§\ref{sec:portage-flink-spark}).

\paragraph{Méthode.} Nous avons suivi la même progression que la thèse de
référence, dans le même ordre : comprendre la théorie avant d'écrire du
code (chapitre~\ref{ch:algorithmes}), décrire le modèle d'exécution cible
avant de l'utiliser (chapitre~\ref{ch:spark}), implémenter opérateur par
opérateur en partant d'une version simple avant d'ajouter des raffinements
(chapitre~\ref{ch:implementation}), et enfin valider la \emph{qualité} du
clustering avant de mesurer sa \emph{scalabilité}
(chapitre~\ref{ch:evaluation}) --- mesurer un speedup sur un algorithme dont
le résultat serait faux n'aurait aucun sens. Cet ordre n'est pas seulement
une méthode de travail : c'est aussi le plan de ce rapport.

\paragraph{Architecture logicielle.} Le code est organisé en trois couches
strictement séparées, une contrainte de conception que nous avons tenue tout
au long du projet :

\begin{itemize}
  \item \texttt{streamkm.core} --- $k$-means, $k$-means++, distance
    euclidienne au carré. Pur Scala, aucune dépendance.
  \item \texttt{streamkm.coreset} --- l'arbre de coreset et la technique
    merge-and-reduce. Pur Scala, sérialisable, testable sans JVM Spark.
  \item \texttt{streamkm.spark} --- l'orchestration Spark (RDD,
    \texttt{treeReduce}, Structured Streaming). Aucune logique
    algorithmique n'y vit : cette couche ne fait que distribuer des appels
    aux deux précédentes.
\end{itemize}

Cette séparation permet de tester et de faire du micro-benchmarking sur les
deux premières couches indépendamment de tout cluster Spark
(chapitre~\ref{ch:evaluation}, §\ref{sec:eval-scalabilite}), et elle rend
visible, dans le code même, la frontière entre « ce qui est StreamKM++ » et
« ce qui est Spark ».
